\section{The registry of approach families}
\label{sec:reg}

The search was run as a portfolio, several mechanisms alive at once and kept independent
in the early rounds so that they would not converge on the same attractive reduction, with
a written registry of what each had achieved and where it was blocked. A route was marked
blocked for a band only together with the measurement that blocked it. Both the successes
and the failures are recorded here, because the failures are what justify the shape of the
final method: they are the reason there is a subdivided band and a $k$-series band at all.

\emph{A caveat on every timing in this section.} The per-offset and per-block times quoted
below were measured in the first round, on a machine carrying four to ten other agents'
Julia processes. They are sound for ordering the families against each other, which is what
they are used for here, and they are \emph{not} the deliverable's timings: those were
re-measured on a quiet machine and are in Sections~\ref{sec:bch} and~\ref{sec:bch:offset}.

\begin{center}\footnotesize
\begin{tabular}{@{}cL{0.24\textwidth}L{0.16\textwidth}L{0.40\textwidth}@{}}
\toprule
& family & status & the number that decided it\\
\midrule
(a) & Cartesian Taylor about the centre & alive, cross-check &
$3\times10^{-16}$--$1.5\times10^{-14}$ from $3$ cells; $\sqrt2\rho<1$ fails at $2$\\
(b) & spherical addition theorem & \textbf{won}, whole box &
$8.6\times10^{-15}$ max-norm over $427$ references\\
(b$'$) & the same on eight octants & \textbf{won}, near band &
the $12$ two-cell offsets at $4.2\times10^{-15}$\\
(c) & one exact axis, two expanded & alive, not used &
$\tau_\xi = 1$ \emph{exactly} at two cells\\
(d) & subdivision of the box & \textbf{won} as (b$'$); reference &
$\rho_j$ falls like $1/(m(n-1)+1)$; the needle needs $\sim\!200$ boxes\\
(e) & spectral / Ewald & blocked, all bands &
$P \approx 1651$ against $256$, i.e.\ $768$~GB, at $f = 1+0.1\ii$\\
(f) & interpolation of the far tensor & alive, marginal &
pays only above $2.2\times10^{3}$ flops/offset; ceiling $9.6\times$\\
(g) & derivative-free coefficients & alive, cross-check &
$4.8\times10^{-15}$ over $290$ rows; agrees with (b) to $1.1\times10^{-14}$\\
(k) & the moment $k$-series & \textbf{won}, slender needle &
$67$ offsets; $\varepsilon\Lambda = 2.3\times10^{-13}$ forces $128$ bits\\
\midrule
\multicolumn{4}{@{}l}{\emph{unequal cells} (Section~\ref{sec:xs}; measured in the
cross-scale round, load $3$--$12$)}\\
(u$_1$) & the whole trapezoid box, Theorem~A$'$ & \textbf{won}, far field &
$98.8\,\%$ of the $\lambda/2$ block; $2.4\times10^{-15}$ over $244$ records\\
(u$_2$) & the gcd average as a box sum over the fine block & \textbf{won}, near band &
$\Lambda = 1.0$--$2.3$; $658$~s naive against $1.4$~s as a box sum\\
(u$_3$) & the twenty-seven-piece split & built, measured, \textbf{dropped} &
$0$ of $2213$ near-band offsets converge at ratio $16$; $5.4$--$8.9\times$ slower
where it does\\
(u$_4$) & the $k$-series on thirty-six unequal panels & fallback, off the lattice &
$103$ records at $10^{-17}$; $6$--$21$~s per offset; $0$ offsets of either block\\
\bottomrule
\end{tabular}
\end{center}

\subsection{(a) Cartesian Taylor expansion about the centre separation}
\label{sec:reg:a}

\textbf{Mechanism.} $\partial_a\partial_b g(\mathbf{R}+\boldsymbol{\delta}) =
\sum_{\boldsymbol{\alpha}}(\partial^{\boldsymbol{\alpha}+\mathbf{e}_a+\mathbf{e}_b}
g)(\mathbf{R})\boldsymbol{\delta}^{\boldsymbol{\alpha}}/\boldsymbol{\alpha}!$, integrated
term by term against the exact moments \eqref{eq:vol:mu}; only even
$\boldsymbol{\alpha}$ survive.

\textbf{Outcome: works from three cells outward, with a proven bound.} The sign and
normalization were pinned against the $220$-bit reference to $8\times10^{-26}$ at
$(5,1,0)$ and $10^{-40}$ at $(8,8,0)$. Over $185$ offsets the measured accuracy is
$3\times10^{-16}$ to $1.5\times10^{-14}$ with $-0.11$ to $+1.84$ digits lost, and the
truncation degree saturates beyond about sixteen cells at $p = 8$, $12$, $16$--$18$ at
$\lambda/32$, $\lambda/8$, $\lambda/4$, giving a constant far cost of $6408$ flops and
$7.7~\mu$s per offset, $34$ to $185$ times faster than the rule it replaces. The
remainder bound was proved with explicit constants --- the exact ball min-modulus
$|\mathbf{R}|-\rho$ for $\rho \le |\mathbf{R}|/2$ and
$\sqrt{|\mathbf{R}|^2/2-\rho^2}$ beyond, a Helmholtz alternative for the diagonal, and a
tail sharpened by the moment weights --- and is $2.0$ to $7.7$ times pessimistic in $p$.

\textbf{Where it stopped.} Its convergence criterion is the \emph{complex-ball} one,
$\sqrt{2}\,r_{\mathrm{d}} < |\mathbf{R}|$, which fails at the six nearest non-touching
axis offsets of a cubic grid. Remark~\ref{rem:bnd:two} settles what that means: the series
itself converges for $\rho < 1$ and it is only the Cauchy bound that is unavailable, but
the practical consequence stands, since at $(2,0,0)$ the best attainable is
$7.7\times10^{-14}$ in \code{Float64} at $p = 100$ ($1.4\times10^{6}$ derivatives,
$9.1\times10^{5}$ sum terms) and the derivatives overflow beyond $p \approx 100$. A
$2^3$ subdivision restores it, at $5.5\times10^{-14}$ with $p = 36$ and
$1.7\times10^{7}$ flops. The slender cell breaks it along the short axis and not
marginally: the criterion demands $n \ge 33$ there, and at $n = 16$ the ratio is
$\tau = 2.00$ and the expansion diverges. An attempt to recover the anisotropy by
replacing the ball with a Cartesian polydisc was made and recorded as not useful: the tail
needs $s_i/r_i < 1$ on every axis, so $r_\perp > s$ always, and for a cubic cell at
$n = 4$ the best achievable $\max_i s_i/r_i$ is $0.69$--$0.77$ against the isotropic
$\tau = 0.61$.

\subsection{(b) Spherical addition theorem, whole box and octant split}
\label{sec:reg:b}

\textbf{Mechanism.} Sections~\ref{sec:exp:add}--\ref{sec:exp:refl}.

\textbf{Outcome: this is the method.} Sign $+$ with no extra factor, bit exact against the
reference at $(5,1,0)$. Against the $220$-bit reference over $163$ tensors from three
cells outward, all four shapes, all three directions and both frequencies, the whole-box
form is accurate to $4.2\times10^{-15}$ with at most $1.3$ digits lost and is rounding
limited (at the same $L$ in \code{BigFloat} the error is $10^{-19}$ to $10^{-22}$);
symmetries and $f$-homogeneity are bit exact. Cost $391$/$1250$/$2731$/$4972$~ns at
$L = 10/20/30/40$ with no allocation, $96.5\,\%$ of the offsets of a $128^3$ block at
$\lambda/32$ needing $L \le 8$, and $2.37$~s for that whole far field single-threaded
with the Miller repair ($0.59$~s without it). The conditioning is confined to the
geometry table: the $l$-sum cancels by at most $8.4$ at $\lambda/32$ and $67$ at
$\lambda/4$, the $n$-sum by $3$ to $4$, while the exact-rational contraction cancels by
$5.8\times10^{8}$ and is therefore not done in floating point at all. The octant split
closes the near band: the twelve two-cell offsets reach $4.2\times10^{-15}$ (worst per
entry $4.3\times10^{-15}$) at both frequencies with sub-box cancellation $\le 1.55$.

\textbf{Where it stopped, and what closed it.} Two gaps, both truncation and neither
conditioning. The twelve offsets with $|\mathbf{n}| \le 2.56$ need $L = 94$ for
$6.6\times10^{-14}$ and do not reach $10^{-15}$ by $L = 100$; the octant split solves
them (Section~\ref{sec:join:twelve}). The slender needle $(0,0,n)$ with $n < 23$ has
$\rho > 1$ and \emph{diverges}; at $n = 24$ the error is $1.3\times10^{-4}$ at
$L = 40$, at $n = 32$ it is $1.5\times10^{-9}$ and only at $n = 64$ does it reach
$3.4\times10^{-16}$. Neither the whole box nor the octant split has a certificate on the
$67$ offsets of Section~\ref{sec:join:slender}, and those go to the $k$-series.

\textbf{Added here.} For a cubic cell the entire $l = 2$ shell vanishes identically
(Remark~\ref{rem:exp:l2}), so the surviving shells are $\{0,4,6,8,\dots\}$; this was not
noticed in the round that produced the term-count tables and is an unclaimed saving on
the shapes Gila uses most.

\subsection{(c) One exact axis, the other two expanded}
\label{sec:reg:c}

\textbf{Mechanism as proposed.} Do the integral along one axis in closed form with the
other two coordinates frozen, expand only in the remaining two, so that the ratio becomes
two dimensional --- the natural tool for a slender cell whose isotropic $\rho$ is
dominated by its long edges.

\textbf{Outcome: the proposal does not exist; an inversion of it works.} With the
transverse coordinates frozen, the line integral of $e^{\ii k\varrho}\varrho^{-l}$ is an
incomplete-Hankel object for every $l \ge 1$, so there is no closed form to have. Reversing
the roles does work: $r^2 = \varrho(t)^2 + \xi$ with $\xi = 2R_2y+y^2+2R_3z+z^2$ a
polynomial in the two transverse coordinates, so $F(r) = H(r^2)$ is expanded in $\xi$
(the operator $D = (1/2r)\dd/\dd r$ acts in closed form on $e^{\ii kr}r^{-l}$), the
transverse moments are exact per level, and the along-$\mathbf{R}$ axis carries a Taylor
series whose coefficients come from the closed system
$G_l' = \ii k(R_1+t)G_{l+1} - l(R_1+t)G_{l+2}$. Accuracy $\le 1.3\times10^{-13}$ per
entry from three cells outward, all directions, $\lambda/32$ and $\lambda/4$, $f = 1$ and
$1+0.1\ii$, with $\le 1.07$ digits lost; $(M^*,Q^*,\mu s) = (25,28,142)$ at three cells,
$(13,22,53)$ at four, $(7,14,18)$ at eight --- comparable with the Taylor route. On the
slender cell it is the best of the Cartesian family: $9~\mu$s at $9.4\times10^{-14}$ for
$(16,0,0)$ and $150~\mu$s for the needle at $(0,0,32)$ against $1.2$~ms.

\textbf{Where it stopped.} At two cells, by an identity rather than by slow convergence:
$\tau_\xi = 1$ \emph{exactly}, even with two exact axes, because the nearest point of the
box is at distance $s$ and the transverse excursion is also $s$. Only subdivision moves
that. The needle at $n \le 16$ is blocked as well ($\tau_\xi = 1.13$ at $n = 16$). And its
$Q$ rule is not proven: the bound with $e^{1.5k R_0\theta}$ is valid but $1.8$ times loose
at $kR = 22$, while the $M$ rule was never violated in $110$ cases.

\subsection{(d) Subdivision of the difference box}
\label{sec:reg:d}

\textbf{Mechanism.} Section~\ref{sec:exp:sub}, with any of (a), (b), (g) as the per-sub-box
engine.

\textbf{Outcome: it supplies the near band, and a third independent reference.} An
independent \code{BigFloat} Gauss--Legendre \emph{volume} rule written for this family
matches the face-pair \code{pairKer} reference to $10^{-63}$--$10^{-65}$ on the cubes and
$1.4\times10^{-62}$ on the slender cell --- the third confirmation of
Theorem~\ref{thm:vol:main} --- and it is the only usable reference on the slender cell,
where \code{pairKer} stalls beyond fifteen minutes per offset. Its own remainder bound
locates the singularities exactly on the complex line at $|t| = |\mathbf{X}|/
|\boldsymbol{\delta}'|$ (the same fact as Proposition~\ref{prop:bnd:rho}) and is never
violated over $364$ rows, with a median looseness of $2.29$ in degree ($1.75$ to $7.74$).
The per-sub-box ratio falls like $1/(m(n-1)+1)$ and not like $1/m$: at $n = 2$,
$m = 1,2,3,4,6,8$ give $0.866$, $0.522$, $0.433$, $0.333$, $0.243$, $0.190$. The optimum
is $m = 2$ or $3$ at two cells and $m = 1$ from three cells outwards; $m = 8$ at two cells
costs $3.5$ times more than $m = 2$, and non-uniform splits gain only $1.00$--$1.09$ on
cubes ($1.5$ on the slender cell). Cancellation between sub-boxes is $1.0$ to $4.0$.

\textbf{Where it stopped.} Two things. Its bound is a certificate and not a selector: a
degree factor of $2.29$ is a flop factor of about twelve, since work grows like $p^3$.
And it does not rescue the needle at an acceptable price: $(0,0,2)$ needs about two
hundred boxes for $\rho_j < 0.5$ ($10\times10\times2$ graded, $p = 25$,
$2.7\times10^{7}$ flops) and $(0,0,8)$ needs $8\times8\times1$ at $p = 21$,
$3.1\times10^{6}$ flops. It converges there; the cost is the objection. One trap is worth
recording because it is invisible in the ratio tables: \emph{odd $m$ is unsafe for an
offset with a zero component}, because a sub-box centred on the zero axis has
$\rho_j = 5.66$ against $0.99$ for the even split at the slender $(0,0,2)$ ---
the same conclusion Lemma~\ref{lem:exp:straddle} reaches from the weight.

\subsection{(e) Spectral and Ewald splitting}
\label{sec:reg:e}

\textbf{Mechanism.} The Fourier transform of $g$ is $1/(|\boldsymbol{\xi}|^2-k^2)$ and
cell averaging is a product of $\mathrm{sinc}^2$ factors, so the whole Toeplitz kernel
might be obtained for all offsets at once by one FFT of an analytic spectral kernel, with
the near field corrected in real space.

\textbf{Outcome: blocked in every band, on structure and not on effort.} The spectrum
$\widehat{T}_{ab} = V_{\mathrm{t}}\prod_i\mathrm{sinc}^2(\xi_is_i/2)\,
(\delta_{ab}k^2-\xi_a\xi_b)/(f^2(|\boldsymbol{\xi}|^2-k^2))$ was verified against real
space to $3.4\times10^{-13}$, so the transform is right. The plain alias sum needs
$M \approx 4\times10^{11}$ images per axis. Ewald damping fixes the aliasing ($M = 2$ at
$\eta s \le 2.5$), but the pole at $|\boldsymbol{\xi}| = k$ is fatal for a different
reason: the inverse DFT of the periodized spectrum is the $P$-periodization of a kernel
that does not decay, with error $e^{-\operatorname{Im}(k)Ps}$. At real frequency that is
an $O(1)$ error at every $P$; at $f = 1+0.1\ii$ and $N = 128$ it needs $P \approx 1651$
against $256$, i.e.\ $768$~GB. And even if the pole were handled, the cost is about
$10^{11}$ flops against $5\times10^{9}$ for the direct route. Reopen only for a mechanism
that removes the pole, not for a better damping.

\subsection{(f) Interpolation of the smooth far tensor}
\label{sec:reg:f}

\textbf{Mechanism.} The tensor entries are analytic in $\mathbf{R}$ away from the origin;
evaluate on a coarse sub-lattice and interpolate.

\textbf{Outcome: alive, marginal, and not needed.} The analyticity radius is exactly the
cell gap. At $\lambda/32$ the Chebyshev degree for $10^{-13}$ is $18$ for a patch of
$H = 32$ cells and $26$ for $H = 64$, and the fraction of offsets that must still be
computed directly is $0.10$ over a $128^3$ grid, so the ceiling is a factor $9.6$. But the
interpolation arithmetic is about $2.1\times10^{3}$ flops per offset, so it pays only if
the direct method costs more than $2.2\times10^{3}$ flops per offset --- which the
whole-box multipole sum at $L \le 8$ does not. Equispaced nodes never reach $10^{-13}$
(Runge), and there is nothing at $\lambda/8$, $\lambda/4$, or in the two-to-eight-cell
band. Kept open as a pure constant-factor option if the per-offset cost of the chosen
method ever rises.

\subsection{(g) Derivative-free Taylor coefficients}
\label{sec:reg:g}

\textbf{Mechanism.} Generate the Taylor coefficients $c_{\boldsymbol{\alpha}}$ of
$g(\mathbf{R}+\boldsymbol{\delta})$ directly from the differential equation the kernel
satisfies, instead of differentiating it.

\textbf{Outcome: the strongest of the Cartesian routes, and the independent cross-check of
(b).} The working route is the pair recurrence on $G = g(\mathbf{R}+\boldsymbol{\delta})$
and $S = |\mathbf{R}+\boldsymbol{\delta}|\,G$: $\nabla S = \ii k(\mathbf{R}+
\boldsymbol{\delta})G$ and $P\nabla G = (\mathbf{R}+\boldsymbol{\delta})(\ii kS - G)$ with
$P = |\mathbf{R}+\boldsymbol{\delta}|^2$, a two-level recurrence seeded by
$a_0 = e^{\ii k\varrho}/(4\pi f^2\varrho)$, $b_0 = e^{\ii k\varrho}/(4\pi f^2)$, and the
assembly
\[
T_{ab} = \frac{1}{V_{\mathrm{t}}}\sum_{\boldsymbol{\beta}\text{ even}}
\mu_{\boldsymbol{\beta}}\Bigl[(\beta_a+1+\delta_{ab})(\beta_b+1)\,
c_{\boldsymbol{\beta}+\mathbf{e}_a+\mathbf{e}_b} + \delta_{ab}k^2c_{\boldsymbol{\beta}}
\Bigr],
\]
which equals \code{srfSum!}'s output with no extra sign or factor. Re-derived
independently and confirmed to $10^{-76}$ at $256$ bits, with the trace identity
$\operatorname{tr}T_q = 2k^2\langle g\rangle_q$ holding to $2\times10^{-58}$. Over $290$
(case, frequency) rows against the $220$-bit reference, and $200$ of them inside Gila's
operating range, the worst tensor error is $4.8\times10^{-15}$ and the worst per entry
$8.5\times10^{-15}$, with $q^* \le 14$ far out. Cost $17.6$--$18.6$~ns per coefficient
and a $32^3$ block in $1.85$~s against $23.3$~s, a $12.6$-fold block speedup.

\textbf{Where it stopped, with numbers.} Three failure modes, all shared with (a) because
it is the same expansion: $\rho \ge 1$ on the slender needle diverges, with the boundary
at $n = 22.6$ exactly and no $(0,0,n)$ with $n \le 39$ reaching $10^{-13}$ within
$q \le 40$; $k\cdot 2s \ge 10$ (cells larger than $\lambda/2$, outside Gila) costs
$10^{0.2k\cdot2s}$; and $\rho \to 1$ at $(2,0,0)$ needs $q \approx 100$ where the
cancellation is $1.8\times10^{4}$, so the error curve turns around at
$1.1\times10^{-13}$. Its imaginary part is limited by the formulation and not by the
recurrence: the tensor is one complex sum, so the relative error of
$\operatorname{Im}T_{ab}$ is $10^{-13}$ to $4\times10^{-12}$ whenever
$|\operatorname{Im}|/\max|T|$ falls below $10^{-3}$ --- which is exactly the defect the
$y_l$/$j_l$ split of Section~\ref{sec:exp:reim} does not have. Four sub-routes were
measured and blocked: the Helmholtz march from exact Cauchy data ($3.5$ digits lost at
order $30$); the addition theorem expanded into Cartesian \emph{monomials} ($3.0$ digits
at order $16$, amplification $10^{0.33n}$, from the $(1+\sqrt2)^l$ growth of the Legendre
coefficients in the monomial basis --- note this is \emph{not} the $(l,m)$-organized
route of Section~\ref{sec:exp}, which is where the same expansion becomes well
conditioned); a one-dimensional ODE with angular recovery (condition $10^{0.33n}$,
$65$--$125$ times more expensive); and automatic differentiation (correct, $37$--$94$
times slower). Two further measurements are worth carrying: order $50$--$60$ at
$kR = 30$ loses $6$--$9$ digits, so the order must be set by $\rho$ and never pushed; and
the flat error floor at large $kR$ is the seed $e^{\ii k\varrho}$, removable for twelve
flops by a double-double phase, which recovers $2.5$ of the $3.35$ digits lost at
$kR = 3000$ --- worth doing for the $h_0$ seed of (b) as well, since $kR$ reaches about
$350$ on a $128^3$ block at $\lambda/4$.

\subsection{\texorpdfstring{(k) The moment $k$-series carried outward}{(k) The moment k-series carried outward}}
\label{sec:reg:k}

\textbf{Mechanism.} The near-field series of the companion document,
$\int_A\int_B g = \frac{1}{4\pi f^2}\sum_n (\ii k)^nI_{n-1}/n!$, applied to separated
pairs.

\textbf{Outcome: correct everywhere, affordable only where $kD$ and $\Lambda$ are both
small.} Its bound (Theorem~\ref{thm:bnd:kser}) is tight to $\pm3$ terms over $233$ rows,
which makes it the sharpest bound in this document. The obstruction is the pair of
constants in front of it. The assembly amplification is $\Lambda = 68$ to $1405$ at
$\lambda/32$, $15$ to $24$ at $\lambda/4$ and $2.2\times10^{3}$ to $8.1\times10^{4}$ for
the slender cell, so the \code{Float64} band edges for $10^{-13}$ on the assembled tensor
are four cells at $\lambda/32$, three at $\lambda/8$, two at $\lambda/4$ and
\emph{none at all} for the slender cell, where $\varepsilon\Lambda = 4.8\times10^{-13}$
already at two cells. In $128$-bit arithmetic the error is $10^{-31}$ to $10^{-28}$
everywhere. And the cost is prohibitive as a general method: the moment ladder scales as
$m_{\max}^{2.6\dots5.1}$ and is re-run per $m$ on three-dimensional pairs, giving $32$~ms
per offset at two cells at $\lambda/32$ (against $13.9$~ms for the rule it would replace)
and $671$~ms at thirty-two cells (against $0.55$~ms).

\textbf{Where it stopped, and the merge that failed.} The idea that would have unified the
near and far fields is the shifted exponential, $e^{\ii kr} = e^{\ii kR_0}
e^{\ii k(r-R_0)}$ with $R_0 = |\mathbf{R}|$: the outer series then needs only $11$ terms
at $\lambda/32$ and $19$--$21$ at $\lambda/4$ with cancellation $\le 1.5$ at every
separation. It fails on the centred moments $\int\int w\,(r-R_0)^n$, which cannot be
generated stably: the naive binomial loses $n\log_{10}(2.3N)$ digits; the divergence
recurrence subtracts $nR_0/d$ per step; and the route through
$q = (2\mathbf{R}\cdot\boldsymbol{\delta}+|\boldsymbol{\delta}|^2)/R_0^2$ converges only
when $(2N+3)/N^2 < 1$, that is $N \ge 4$ on an axis --- which is family (a)'s own radius,
so the merge is circular. Blocked as a merge; the series stays what it is, the tool for
the touching shell and for the slender needle, where $kD_{\max} \le 0.60$ and
$N \le 16$.

\textbf{Side findings.} Two, both about the code being replaced. Gila's fixed rule on the
slender cell at $(0,0,8)$ has a relative tensor error of $1.79\times10^{-3}$ (order-$5$
rule; $2.8\times10^{-5}$ on the worst face pair), and at $(0,0,20)$ and $(0,0,24)$ it is
$4.5\times10^{-6}$ and $1.9\times10^{-6}$, because \code{quadOrd} keys the order on the
separation in cells and not on the physical geometry. And the \code{Float64} moment
ladder itself loses $5.5$ digits at $m_{\max} = 48$.

\subsection{Three disagreements between rounds, and how they were resolved}
\label{sec:reg:disc}

\begin{enumerate}[nosep,leftmargin=1.6em]
\item \emph{Does the Cartesian series converge at the first axis offset?} One round
recorded $(2,0,0)$ as divergent and asymptotic; another recorded it as convergent with an
effective ratio of $0.73$. Proposition~\ref{prop:bnd:rho} and the measurement in
Remark~\ref{rem:bnd:two} settle it: convergent, with an observed asymptotic ratio of
$0.887$ tending to $\rho = 0.866$, and no Cauchy \emph{bound} available. The practical
conclusion --- hand the offset to the octant split --- is unchanged, but the reason is
different, and the difference matters for the slender needle, where $\rho > 1$ and the
divergence is real.
\item \emph{The sub-box term count.} The bound \eqref{eq:bnd:sub} selects
$L_j \le 56$ at the twelve two-cell offsets when evaluated by \code{bands.jl} and
$L_j \le 48$ when evaluated in the round that introduced it, from the same written
formula and the same geometry. The discrepancy is about a factor $100$ in the tail, does
not move any band edge, and is recorded unresolved in Remark~\ref{rem:join:disc}; the
larger value is used here because it is what the stated formula produces.
\item \emph{The truncation rule for the Taylor routes.} One round proposed
$q^* \approx 17/\log_{10}(1/\rho)$ and the audit measured $q^* \approx 7 +
7.1/\log_{10}(1/\rho)$, the first being $3.35$ times too costly and a pure
$9.7/\log_{10}(1/\rho)$ undershooting by two at $\rho \le 0.06$, where ninety per cent of
a block's offsets sit. The audit's rule is the one to use; it is a fit and not a bound,
which is a reason to prefer the multipole form, whose $L$ comes from
Theorem~\ref{thm:bnd2:reg}.
\end{enumerate}

\subsection{The second round: what was rebuilt after the families closed}
\label{sec:reg:r2}

The first round settled the mathematics and produced no library. The second round produced
one and then attacked it. Six passes, in order, each with the number that made it necessary.

\begin{description}[leftmargin=1.4em]
\item[The theory pass.] Wrote Sections~\ref{sec:vol}--\ref{sec:join} from the round-1
reports and checked every bound in \code{BigFloat} against the quantity it bounds. It found
that the constant $17$ of Theorem~\ref{thm:bnd:rem} is loose by $32$ to $34\,000$; that the
end-to-end bound over-states by $1.3\times10^{3}$ to $3.0\times10^{7}$; that for a cubic
cell the entire $l = 2$ shell vanishes (Remark~\ref{rem:exp:l2}); and that the band edges
follow from the bounds alone (Section~\ref{sec:join:bands}) rather than from experiment.

\item[The unification pass.] Built \code{farfield.jl}: the three routes, the exact geometry
table, the selector, the caches, threading. Accuracy $8.6\times10^{-15}$ max-norm over the
$403$ references then available; the $128^3$ far field in $2.05$~s single-threaded. It also
found four defects in the round-1 sources it superseded --- a tail bound that was a partial
sum and therefore vacuous at its own top shell, a geometry table that accepted
\code{Float64} edge lengths and became noise beyond $L \approx 26$, a Bessel seed losing
$7.7\times10^{-14}$ of phase at $k\mathbf{R} = 348$, and a per-thread workspace sized by
\code{nthreads()} instead of \code{maxthreadid()}.

\item[The sharper bound.] Replaced the two lossy steps of Theorem~\ref{thm:bnd:rem} by exact
identities (Section~\ref{sec:bnd2}): a median gain of $3.2\times10^{3}$ in the bound, of
$1.47$ in per-offset cost, and --- the point --- a bound on the truncation the code actually
performs rather than on a different reorganization of the same series.

\item[The preparation pass.] Integrated Theorem~A/B into the selector, moved the tolerance
to $10^{-14}$ for the reason in Section~\ref{sec:join:rule}, rebuilt the geometry table at
$256$ bits (bit-for-bit identical to the $1024$-bit table over all $321\,373$ stored numbers
of four shapes) and sized it lazily in $L$, cutting a cold $32^3$ build from $828$~s to
$202$~s and one cold far offset from $828$~s to $1.3$~s. It also found that the
$n$-truncation cut was being evaluated at a radius where the expansion does not converge,
that the exact-zero threshold inside the contraction did not track the build precision, and
that the selection arithmetic must be done in \code{Float64} whatever the working type,
because $(2l+1)!!$ overflows a \code{Float32} at $l = 29$.

\item[The adversarial audit.] Twelve cell shapes with aspect ratios $1$ to $1020$, six
shape-adapted offsets each, four frequencies: $288$ tensors against an independent $220$-bit
graded volume rule whose own convergence is a median $6.0\times10^{-50}$; $60$ further
random cases for the bound; $40$ boundary offsets; $21$ frequency cases including $f = -1$,
$f = 1-0.1\ii$, $f = 10$ and $f = 1+3\ii$; determinism at $1$, $4$ and $12$ threads. It
produced the ten defects of Section~\ref{sec:def}, of which D9 --- the $k$-series of $j_l$
capped at a compile-time constant --- destroyed the answer completely
($9.6\times10^{-3}$ per entry) as soon as $|k| r_{\mathrm{d}}$ exceeded about $3.2$.

\item[The merge.] Fixed D9 by Remark~\ref{rem:bnd:adapt}, confirmed the audit's other two
accuracy defects fixed, removed a cost model that was computed on every offset and read by
nobody ($28$--$31\,\%$ of the whole far-field fill), threaded the $k$-series route
($633$~s $\to$ $128$~s on a cold slender block, bitwise identical), and closed the five API
gaps. It also found, while testing positivity, that the indefiniteness two earlier passes
had attributed to the discretization is Gila's order-$9$ contact rule
(Section~\ref{sec:ver:pos}).
\end{description}

\subsection{The cross-scale round: unequal cells}
\label{sec:reg:x}

The third round extended the library to pairs of cells with different edges
(Section~\ref{sec:xs}). Five agents ran in parallel from one brief --- the theory, the
table and routes, Gila's geometry and accuracy, the $220$-bit references, and the
near-field item --- followed by an integration pass and an adversarial one. Four families
were alive; two won, one is the fallback, one was dropped.

\begin{description}[leftmargin=1.4em]
\item[(u$_1$) The whole trapezoid box.] \textbf{Mechanism.} The equal-cell table with the
trapezoid moments $\mu_n(a,b) = \mu_n(b)-\mu_n(a)$ in place of the triangle's; Theorem~A
with the ramp-only flat measure and point masses at $\pm a_i$, $\pm b_i$
(Theorem~\ref{thm:xs:reg}). \textbf{Outcome: the far route.} One table per unordered pair,
$3.3$--$5.7$~s cold; $98.8\,\%$ of the realistic $\lambda/2$ block at $0.9$--$12.5~\mu$s
per offset; $2.4\times10^{-15}$ worst over $244$ reference records; the equal-cell tables
bit-identical at $a = 0$. \textbf{Where it stops.} At $\rho_{\mathrm{whl}} = 0.58$--$0.60$,
the same edge as the equal-cell whole box; $16\,908$ offsets of the block lie above it. Its
first results at the $\lambda/4$ and $\lambda/2$ coarse cells were wrong by up to
$6.8\times10^{-10}$, which was the $n$-cut defect D11 and not the table.

\item[(u$_2$) The gcd average as a box sum.] \textbf{Mechanism.} Every Gila pair lives on
the common gcd lattice, so
$G(\mathbf{R};\mathbf{s}_{\mathrm{t}},\mathbf{s}_{\mathrm{s}}) = N_{\mathrm{t}}^{-1}\sum_{jj'}
G(\mathbf{R}+\mathbf{c}_j-\mathbf{c}'_{j'};\mathbf{g},\mathbf{g})$ over integer fine
offsets, evaluated as a weighted sum of entries of the fine equal-cell \code{egoToe}
(Section~\ref{sec:xs:routes}). \textbf{Outcome: the near band.} No new table, no new bound,
all weights $+1/N_{\mathrm{t}}$; cancellation $\Lambda = 1.0$--$2.3$ measured over the
matrix at three frequencies (the a priori fear that it would cancel at $\lambda/2$ was
wrong); $3.4\times10^{-15}$ worst on a $4096$-term sum. \textbf{Where it stopped, once.}
Evaluated naively as $N_{\mathrm{t}}N_{\mathrm{s}}$ separate tensors per offset it costs
$38.9$~ms per offset and $658$~s for the band, $160$ times the far field; as a box sum
$1.4$~s. The formula was right and the first implementation of it was the wrong one.

\item[(u$_3$) The twenty-seven-piece split.] \textbf{Mechanism.} The sub-box route of
Section~\ref{sec:exp:sub} with three pieces per axis, two ramps and a plateau, eight
piece-type tables and Theorem~B with $\beta = 0$ on a plateau axis. \textbf{Outcome:
built in full, dropped.} The all-plateau piece has radius $|\mathbf{a}| = (r-1)\sqrt3g/2$
against a nearest offset of $(r+3)g/2$, a ratio of $0.87$, $1.17$, $1.24$ at $r = 4, 8, 16$:
it diverges exactly where the whole box fails. Where it converges it is $5.4$--$8.9$ times
slower than the whole box and as accurate ($1.5\times10^{-15}$ over $152$ records). It
served as the third route in every cross-check and is kept in
\code{scratch/xwork/xtable/farx.jl}.

\item[(u$_4$) The $k$-series on unequal panels.] \textbf{Mechanism.} Route (iii) with
\code{pairMoments} on the thirty-six face pairs of two different rectangles.
\textbf{Outcome: the correctness fallback.} $10^{-17}$ against $103$ off-lattice records,
$6$--$21$~s per offset at $N = 13$--$44$; reached by $0$ offsets of either block, because
Gila's geometry never leaves the gcd lattice.
\end{description}

\noindent
Two disagreements were resolved on the way. The root's prediction that the gcd sum would
cancel at $\lambda/4$--$\lambda/2$ coarse cells, made to explain whole-box against gcd
disagreements of $5.9\times10^{-14}$--$6.8\times10^{-10}$, was wrong: $\Lambda$ was
$1.1$--$1.6$ at those offsets and the whole box was the inaccurate side (D11). The
root's prediction that the plateau piece of the split would not converge for $r\ge8$ held,
and the measurement extended it to $r = 4$ on two and three axes.
